%%% Sample chapter with typesetting examples

\chapter{Typesetting Guide}

\section{Formatting the thesis}

The thesis text itself is organized hierarchically into chapters and sections, with each chapter starting on a new page. Text is justified. A new paragraph is usually separated by a small vertical space and indentation of the first line. The visual formatting should be consistent throughout the text.

Abbreviations used in the text must be explained at their first occurrence (in parentheses or in a footnote, if a more complex explanation of the term or abbreviation is needed). If there are many abbreviations, a list of abbreviations used should be added, including their explanation and/or references to their definition.

Longer text taken over from another author must be delimited with quotation marks or otherwise marked, and properly cited.

\section{Simple examples}

A narrow (non-breaking) space belongs between a number and a unit: the width of an A4 page is $210\,\rm mm$, which only $5\,\%$ of authors remember. But when the value is used attributively, the space is omitted: a $25\rm mm$ margin, a $95\%$ confidence interval.

We distinguish between different kinds of dashes:
red-black (hyphen),
pages 16--22 (en dash),
$45-44$ (mathematical minus),
and this is --- as you probably expected --- a parenthetical clause set off by em dashes.

Don't hard-code quotation marks with \texttt{``like this''}; use \verb|\enquote{like this}| instead. \enquote{Like this}, it automatically follows the typographic convention of whichever language babel is currently set to, so the same source works unchanged if the thesis language ever changes.

% This paragraph makes the tilde visible
{
\def~{{\tt\char126}}
In several places you need to prevent a line break (in \TeX{} we mark this with a tilde): between a number/letter and the unit or word that must stay with it -- vertex~$v$, in $k$~steps, Fig.~3, pp.~12--15, Ch.~5 -- and generally anywhere a line break would leave the reader \enquote{stumbling}, such as before a short word stranded at the start or end of a line.
}
